%% fig_mfd_concept.tex
%% マクロ基本ダイアグラム（MFD）の概念図
%% Production G(n) [台・km/h] vs Accumulation n [台]
%%
%% パラメータ（正規化単位）:
%%   G(n) = 3n - n²/2 = n(3 - n/2)  （TikZ座標）
%%   → 実際の G(n) = (n/n_j)(1 - n/n_j) × G_max に相当
%%   臨界集積: n* = 3 (TikZ x), G(n*) = 4.5 (TikZ y)
%%   ジャム: n_j = 6 (TikZ x)
%%
%% Usage: %% fig_mfd_concept.tex
%% マクロ基本ダイアグラム（MFD）の概念図
%% Production G(n) [台・km/h] vs Accumulation n [台]
%%
%% パラメータ（正規化単位）:
%%   G(n) = 3n - n²/2 = n(3 - n/2)  （TikZ座標）
%%   → 実際の G(n) = (n/n_j)(1 - n/n_j) × G_max に相当
%%   臨界集積: n* = 3 (TikZ x), G(n*) = 4.5 (TikZ y)
%%   ジャム: n_j = 6 (TikZ x)
%%
%% Usage: %% fig_mfd_concept.tex
%% マクロ基本ダイアグラム（MFD）の概念図
%% Production G(n) [台・km/h] vs Accumulation n [台]
%%
%% パラメータ（正規化単位）:
%%   G(n) = 3n - n²/2 = n(3 - n/2)  （TikZ座標）
%%   → 実際の G(n) = (n/n_j)(1 - n/n_j) × G_max に相当
%%   臨界集積: n* = 3 (TikZ x), G(n*) = 4.5 (TikZ y)
%%   ジャム: n_j = 6 (TikZ x)
%%
%% Usage: %% fig_mfd_concept.tex
%% マクロ基本ダイアグラム（MFD）の概念図
%% Production G(n) [台・km/h] vs Accumulation n [台]
%%
%% パラメータ（正規化単位）:
%%   G(n) = 3n - n²/2 = n(3 - n/2)  （TikZ座標）
%%   → 実際の G(n) = (n/n_j)(1 - n/n_j) × G_max に相当
%%   臨界集積: n* = 3 (TikZ x), G(n*) = 4.5 (TikZ y)
%%   ジャム: n_j = 6 (TikZ x)
%%
%% Usage: \input{assets/assignment/fig_mfd_concept.tex}

\begin{tikzpicture}[
  every node/.style={font=\small},
  xscale=1.0,
  yscale=1.0
]

%% ===== 背景：自由流領域（水色）・過密領域（橙） =====
\fill[cyan!10]  (0, 0) -- plot[domain=0:3, samples=60]
    (\x, {3*\x - 0.5*\x*\x}) -- (3, 0) -- cycle;
\fill[orange!12] (3, 0) -- plot[domain=3:6, samples=60]
    (\x, {3*\x - 0.5*\x*\x}) -- (6, 0) -- cycle;

%% ===== MFD 曲線 =====
\draw[very thick, blue!70!black]
  plot[domain=0:6, samples=100]
  (\x, {3*\x - 0.5*\x*\x});

%% ===== 補助破線：n* と G_max =====
\draw[gray, dashed, thin] (3, 0) -- (3, 4.5);
\draw[gray, dashed, thin] (0, 4.5) -- (3, 4.5);

%% ===== x 軸 =====
\draw[->] (-0.3, 0) -- (6.9, 0)
    node[right] {Accumulation $n$ [台]};

%% x 軸目盛り
\draw (3, 0.08) -- (3, -0.12)
    node[below] {$n^*$};
\draw (6, 0.08) -- (6, -0.12)
    node[below] {$n_{\rm jam}$};
\node[below] at (0, -0.12) {$0$};

%% ===== y 軸 =====
\draw[->] (0, -0.3) -- (0, 5.3)
    node[above, align=center] {Production\\$G(n)$ [台・km/h]};

%% y 軸目盛り
\draw (-0.08, 4.5) -- (0.08, 4.5)
    node[left, xshift=-0.05cm] {$G_{\max}$};
\node[left] at (-0.05, 0) {$0$};

%% ===== G_max 点にマーカー =====
\fill[blue!70!black] (3, 4.5) circle (2.5pt);

%% ===== 領域ラベル =====
\node[align=center, font=\footnotesize, cyan!60!black]
    at (1.5, 1.2) {自由流\\領域};
\node[align=center, font=\footnotesize, orange!70!black]
    at (4.5, 1.2) {過密\\（グリッドロック）};

%% ===== n 方向の注釈 =====
\node[font=\footnotesize, align=center, blue!60!black]
    at (3.2, 4.8) {$n = n^*$（臨界）};

%% ===== 平均速度の参考ライン（原点からの傾きが平均速度に比例） =====
%% 自由流域の代表点 n=1: G(1)=2.5 なので slope=2.5
\draw[gray!80, dotted, thin, ->]
    (0, 0) -- (2.0, 5.0)
    node[above right=-0.1cm, font=\footnotesize, gray!90]
    {$\bar{v} \propto G(n)/n$};

\end{tikzpicture}


\begin{tikzpicture}[
  every node/.style={font=\small},
  xscale=1.0,
  yscale=1.0
]

%% ===== 背景：自由流領域（水色）・過密領域（橙） =====
\fill[cyan!10]  (0, 0) -- plot[domain=0:3, samples=60]
    (\x, {3*\x - 0.5*\x*\x}) -- (3, 0) -- cycle;
\fill[orange!12] (3, 0) -- plot[domain=3:6, samples=60]
    (\x, {3*\x - 0.5*\x*\x}) -- (6, 0) -- cycle;

%% ===== MFD 曲線 =====
\draw[very thick, blue!70!black]
  plot[domain=0:6, samples=100]
  (\x, {3*\x - 0.5*\x*\x});

%% ===== 補助破線：n* と G_max =====
\draw[gray, dashed, thin] (3, 0) -- (3, 4.5);
\draw[gray, dashed, thin] (0, 4.5) -- (3, 4.5);

%% ===== x 軸 =====
\draw[->] (-0.3, 0) -- (6.9, 0)
    node[right] {Accumulation $n$ [台]};

%% x 軸目盛り
\draw (3, 0.08) -- (3, -0.12)
    node[below] {$n^*$};
\draw (6, 0.08) -- (6, -0.12)
    node[below] {$n_{\rm jam}$};
\node[below] at (0, -0.12) {$0$};

%% ===== y 軸 =====
\draw[->] (0, -0.3) -- (0, 5.3)
    node[above, align=center] {Production\\$G(n)$ [台・km/h]};

%% y 軸目盛り
\draw (-0.08, 4.5) -- (0.08, 4.5)
    node[left, xshift=-0.05cm] {$G_{\max}$};
\node[left] at (-0.05, 0) {$0$};

%% ===== G_max 点にマーカー =====
\fill[blue!70!black] (3, 4.5) circle (2.5pt);

%% ===== 領域ラベル =====
\node[align=center, font=\footnotesize, cyan!60!black]
    at (1.5, 1.2) {自由流\\領域};
\node[align=center, font=\footnotesize, orange!70!black]
    at (4.5, 1.2) {過密\\（グリッドロック）};

%% ===== n 方向の注釈 =====
\node[font=\footnotesize, align=center, blue!60!black]
    at (3.2, 4.8) {$n = n^*$（臨界）};

%% ===== 平均速度の参考ライン（原点からの傾きが平均速度に比例） =====
%% 自由流域の代表点 n=1: G(1)=2.5 なので slope=2.5
\draw[gray!80, dotted, thin, ->]
    (0, 0) -- (2.0, 5.0)
    node[above right=-0.1cm, font=\footnotesize, gray!90]
    {$\bar{v} \propto G(n)/n$};

\end{tikzpicture}


\begin{tikzpicture}[
  every node/.style={font=\small},
  xscale=1.0,
  yscale=1.0
]

%% ===== 背景：自由流領域（水色）・過密領域（橙） =====
\fill[cyan!10]  (0, 0) -- plot[domain=0:3, samples=60]
    (\x, {3*\x - 0.5*\x*\x}) -- (3, 0) -- cycle;
\fill[orange!12] (3, 0) -- plot[domain=3:6, samples=60]
    (\x, {3*\x - 0.5*\x*\x}) -- (6, 0) -- cycle;

%% ===== MFD 曲線 =====
\draw[very thick, blue!70!black]
  plot[domain=0:6, samples=100]
  (\x, {3*\x - 0.5*\x*\x});

%% ===== 補助破線：n* と G_max =====
\draw[gray, dashed, thin] (3, 0) -- (3, 4.5);
\draw[gray, dashed, thin] (0, 4.5) -- (3, 4.5);

%% ===== x 軸 =====
\draw[->] (-0.3, 0) -- (6.9, 0)
    node[right] {Accumulation $n$ [台]};

%% x 軸目盛り
\draw (3, 0.08) -- (3, -0.12)
    node[below] {$n^*$};
\draw (6, 0.08) -- (6, -0.12)
    node[below] {$n_{\rm jam}$};
\node[below] at (0, -0.12) {$0$};

%% ===== y 軸 =====
\draw[->] (0, -0.3) -- (0, 5.3)
    node[above, align=center] {Production\\$G(n)$ [台・km/h]};

%% y 軸目盛り
\draw (-0.08, 4.5) -- (0.08, 4.5)
    node[left, xshift=-0.05cm] {$G_{\max}$};
\node[left] at (-0.05, 0) {$0$};

%% ===== G_max 点にマーカー =====
\fill[blue!70!black] (3, 4.5) circle (2.5pt);

%% ===== 領域ラベル =====
\node[align=center, font=\footnotesize, cyan!60!black]
    at (1.5, 1.2) {自由流\\領域};
\node[align=center, font=\footnotesize, orange!70!black]
    at (4.5, 1.2) {過密\\（グリッドロック）};

%% ===== n 方向の注釈 =====
\node[font=\footnotesize, align=center, blue!60!black]
    at (3.2, 4.8) {$n = n^*$（臨界）};

%% ===== 平均速度の参考ライン（原点からの傾きが平均速度に比例） =====
%% 自由流域の代表点 n=1: G(1)=2.5 なので slope=2.5
\draw[gray!80, dotted, thin, ->]
    (0, 0) -- (2.0, 5.0)
    node[above right=-0.1cm, font=\footnotesize, gray!90]
    {$\bar{v} \propto G(n)/n$};

\end{tikzpicture}


\begin{tikzpicture}[
  every node/.style={font=\small},
  xscale=1.0,
  yscale=1.0
]

%% ===== 背景：自由流領域（水色）・過密領域（橙） =====
\fill[cyan!10]  (0, 0) -- plot[domain=0:3, samples=60]
    (\x, {3*\x - 0.5*\x*\x}) -- (3, 0) -- cycle;
\fill[orange!12] (3, 0) -- plot[domain=3:6, samples=60]
    (\x, {3*\x - 0.5*\x*\x}) -- (6, 0) -- cycle;

%% ===== MFD 曲線 =====
\draw[very thick, blue!70!black]
  plot[domain=0:6, samples=100]
  (\x, {3*\x - 0.5*\x*\x});

%% ===== 補助破線：n* と G_max =====
\draw[gray, dashed, thin] (3, 0) -- (3, 4.5);
\draw[gray, dashed, thin] (0, 4.5) -- (3, 4.5);

%% ===== x 軸 =====
\draw[->] (-0.3, 0) -- (6.9, 0)
    node[right] {Accumulation $n$ [台]};

%% x 軸目盛り
\draw (3, 0.08) -- (3, -0.12)
    node[below] {$n^*$};
\draw (6, 0.08) -- (6, -0.12)
    node[below] {$n_{\rm jam}$};
\node[below] at (0, -0.12) {$0$};

%% ===== y 軸 =====
\draw[->] (0, -0.3) -- (0, 5.3)
    node[above, align=center] {Production\\$G(n)$ [台・km/h]};

%% y 軸目盛り
\draw (-0.08, 4.5) -- (0.08, 4.5)
    node[left, xshift=-0.05cm] {$G_{\max}$};
\node[left] at (-0.05, 0) {$0$};

%% ===== G_max 点にマーカー =====
\fill[blue!70!black] (3, 4.5) circle (2.5pt);

%% ===== 領域ラベル =====
\node[align=center, font=\footnotesize, cyan!60!black]
    at (1.5, 1.2) {自由流\\領域};
\node[align=center, font=\footnotesize, orange!70!black]
    at (4.5, 1.2) {過密\\（グリッドロック）};

%% ===== n 方向の注釈 =====
\node[font=\footnotesize, align=center, blue!60!black]
    at (3.2, 4.8) {$n = n^*$（臨界）};

%% ===== 平均速度の参考ライン（原点からの傾きが平均速度に比例） =====
%% 自由流域の代表点 n=1: G(1)=2.5 なので slope=2.5
\draw[gray!80, dotted, thin, ->]
    (0, 0) -- (2.0, 5.0)
    node[above right=-0.1cm, font=\footnotesize, gray!90]
    {$\bar{v} \propto G(n)/n$};

\end{tikzpicture}
